\section{Associated Bundles and Covariant Derivative}
\subsection{Associated Bundles}
Let $\pi: P\to M$ be a principal $G$-bundle, $F$ be a smooth manifold with a left $G$-action $\rho:G\to\mathrm{Diff}(F)$. 
\begin{definition}
    The \textbf{associated bundle} of $P$ with fiber $F$ is the quotient space 
    \begin{equation*}
        E := P\times_\rho F := (P\times F)/\sim,
    \end{equation*}
    where the equivalence relation $\sim$ is defined by 
    \begin{equation*}
        (p, f)\sim (pg, g^{-1}f),\quad \forall p\in P, f\in F, g\in G.
    \end{equation*} 
    Denote the equivalence class of $(p, f)$ by $[p, f]$.
    We can define a projection map $\pi_E: E\to M$ by $\pi_E([p, f]) = \pi(p)$. Since 
    \begin{equation*}
        \pi_E([pg, g^{-1}f]) = \pi(pg) = \pi(p) = \pi_E([p, f]),
    \end{equation*}
    $\pi_E$ is well-defined. 
\end{definition}
\begin{example}
    If $P = U\times G$ is a trivial principal $G$-bundle over $U$, 
    then $(U\times G)\times_\rho F$ is isomorphic to $U\times F$. Since 
    \begin{align*}
        (U\times G)\times_\rho F \to U\times F,\qquad
        [(x, g), f] = [(x, e)\cdot g, f] = [(x, e), g\cdot f] \mapsto (x, g\cdot f),
    \end{align*}
    is a well-defined diffeomorphism.
\end{example}
\begin{fact}
    The associated bundle $E:=P\times_\rho F$ with the projection $\pi_E$ gives a fiber bundle over $M$ with fiber $F$.
    % Moreover, if $\{(U_\alpha, \phi_\alpha)\}$ is a local trivialization of $P$ with transition functions $g_{\alpha\beta}: U_\alpha\cap U_\beta\to G$, 
    % then we can construct a local trivialization $\{(U_\alpha, \psi_\alpha)\}$ of $P\times_\rho F$ with transition functions 
    % $\rho\circ g_{\alpha\beta}: U_\alpha\cap U_\beta\to \mathrm{Diff}(F)$.
    Let $\{(U_\alpha, \phi_\alpha)\}$ be a local trivialization of $P$ 
    \begin{equation*}
        \phi_\alpha: \pi^{-1}(U_\alpha) \to U_\alpha\times G, \quad p\mapsto \Big(\pi(p), g_\alpha(p)\Big),
    \end{equation*}
    The definition of $g_\alpha(p)$ is as follows: 
    Let $s_\alpha: U_\alpha\to P$ be a local section defined by $s_\alpha(x) = \phi_\alpha^{-1}(x, e)$. 
    Since the action of $G$ on $P$ is free and transitive on each fiber, for any $p\in \pi^{-1}(U_\alpha)$, 
    there exists a unique $g_\alpha(p)\in G$ such that $p=s_\alpha(\pi(p))g_\alpha(p)$. 
    The transition functions of this local trivialization are given by 
    \begin{equation*}
        g_{\alpha\beta}(x) = g_\alpha(s_\beta(x)).
    \end{equation*}

    Then we can construct a local trivialization $\{(U_\alpha, \psi_\alpha)\}$ of $E$ by 
    \begin{equation*}
        \psi_\alpha: \pi_E^{-1}(U_\alpha) \to U_\alpha\times F, \quad [p, f]\mapsto \Big(\pi(p), \rho(g_\alpha(p))f\Big).
    \end{equation*}
    One can check that $\psi_\alpha$ is well-defined and is a diffeomorphism. 
    The transition functions of this local trivialization are 
    \begin{equation*}
        \rho\circ g_{\alpha\beta}: U_\alpha\cap U_\beta\to \mathrm{Diff}(F).
    \end{equation*}
\end{fact}
\begin{example}
    If $P=M\times G$ is a trivial, then $P\times_\rho F$ is trivial. 
    If the action $\rho$ is trivial, then $P\times_\rho F$ is also trivial.
\end{example}
\begin{example}
    If $F=V$ is a finite-dimensional vector space and $\rho:G\to \mathrm{GL}(V)$ is a linear representation of $G$, 
    then the associated bundle $E=P\times_\rho V$ is a vector bundle over $M$ with fiber $V$. The vector bundle structure is given by
    \begin{equation*}
        [p, v]+[p, w] := [p, v+w], \quad \lambda[p, v] := [p, \lambda v], \quad \forall p\in P, v,w\in V, \lambda\in\mathbb{R}.
    \end{equation*}
    One can check that the vector bundle structure is well-defined.
\end{example}
\begin{example}
    Let $P=\mathrm{Fr}(E)$ be the frame bundle of a rank $m$ vector bundle $E$ over $M$. 
    \begin{itemize}[leftmargin=*]
        \item $V=\mathbb{R}^m$, $\mathrm{GL}(m, \mathbb{R})\curvearrowright\mathbb{R}^m:g\cdot v = gv$, we get $\mathrm{Fr}(E)\times_{\rho_{\mathrm{can}}}\mathbb{R}^m\cong E$. 
        \item $V=\mathbb{R}^{m\times m}$, $\mathrm{GL}(m,\mathbb{R})\curvearrowright\mathbb{R}^{m\times m}:g\cdot A = gAg^{-1}$, we get $\mathrm{End}(E)$. 
        \item $V=\bigwedge^k(\mathbb{R}^m)^*$, $\mathrm{GL}(m,\mathbb{R})\curvearrowright\bigwedge^k(\mathbb{R}^m)^*:g\cdot\alpha(v_1,\ldots,v_k) = \alpha(g^{-1}v_1,\ldots,g^{-1}v_k)$, we get $\bigwedge^k E^*$. 
        \item $V=\mathrm{Sym}^k(\mathbb{R}^m)^*$, $\mathrm{GL}(m,\mathbb{R})\curvearrowright\mathrm{Sym}^k(\mathbb{R}^m)^*:g\cdot \alpha(v_1,\ldots,v_k) = \alpha(g^{-1}v_1,\ldots,g^{-1}v_k)$, we get $\mathrm{Sym}^k E^*$.
        \item $V=\mathbb{R}$, $\mathrm{GL}(m,\mathbb{R})\curvearrowright\mathbb{R}:g\cdot r = (\det g)r$, we get $\det E$.
    \end{itemize}
    % If we take $V=\mathbb{R}^m$ and $\rho_{\mathrm{can}}=\mathrm{id}:\mathrm{GL}(m, \mathbb{R})\to \mathrm{GL}(\mathbb{R}^m)$ is the canonical representation, 
    % then the associated bundle $\mathrm{Fr}(E)\times_{\rho_{\mathrm{can}}}\mathbb{R}^m$ is isomorphic to $E$. 
    % If we take $V=\mathbb{R}^{m\times m}$ and $\mathrm{GL}(m,\mathbb{R})$ acts on $\mathbb{R}^{m\times m}$ by $g\cdot A = gAg^{-1}$,
    % then the associated bundle is isomorphic to $\mathrm{End}(E)$. 
    % If we take $V=\bigwedge^k\mathbb{R}^m$ and $\mathrm{GL}(m,\mathbb{R})$ acts on $\bigwedge^k\mathbb{R}^m$ by $g\cdot (v_1\wedge\cdots\wedge v_k) = gv_1\wedge\cdots\wedge gv_k$, 
    % then the associated bundle is isomorphic to $\bigwedge^k E$.
\end{example}
\begin{example}
    Take $F=G$ and $\rho:G\to \mathrm{Aut}(G):\rho(g)(h) = ghg^{-1}$ the conjugation action, 
    then the associated bundle $P\times_{\mathrm{conj}} G$ is denoted by $\mathrm{Ad}P$.
    Take $V=\mathfrak{g}$ and $\rho=\mathrm{Ad}:G\to \mathrm{GL}(\mathfrak{g})$ the adjoint representation. 
    The associated bundle $\mathrm{ad}P := P\times_{\mathrm{Ad}}\mathfrak{g}$ 
    is called the \textbf{adjoint bundle} of $P$.
\end{example}
Over this section, we mainly focus on the case where $F=V$ and $\rho:G\to \mathrm{GL}(V)$ is a linear representation of $G$. 

Fix $x\in M$, let $P_x$ be the fiber of $P$ over $x$ and $E_x$ be the fiber of $E$ over $x$. 
Then for every $p\in P_x$, there is a canonical way to identify $E_x$ with $V$ by 
\begin{equation*}
    f_p: V\to E_x, \quad f_p(v) = [p, v].
\end{equation*}
\begin{lemma}
    The map $f_p$ is a linear isomorphism.
\end{lemma}
\begin{proof}
    Injectivity: If $[p, v]=[p, w]$, then there exists $g\in G$ such that $pg=p$ and $g^{-1}v=w$. 
    Since the action of $G$ on $P$ is free, we have $g=e$ and thus $v=w$. 

    Surjectivity: For any $[q,w]\in E_x$, $q\in P_x$, so there exists $g\in G$ such that $q=pg$. 
    Then we have $[q,w] = [pg, w] = [p, gw] = f_p(gw)$. 
\end{proof}
\begin{lemma}\label{lem:G-equivariance-of-fp}
    Let $f_p: V\to E_x$ be the map defined above. 
    For any $g\in G$, $f_{pg} = f_p\circ \rho(g)$.
\end{lemma}
\begin{proof}
    For any $v\in V$, we have 
    \begin{equation*}
        f_{pg}(v) = [pg, v] = [p, gv] = f_p(gv) = f_p\circ \rho(g)(v). \qedhere
    \end{equation*}
\end{proof}

% \subsubsection*{The Local Trivialization of an Associated Bundle}
% Let $\{(U_\alpha, \phi_\alpha)\}$ be a local trivialization of $P$ 
% \begin{equation*}
%     \phi_\alpha: \pi^{-1}(U_\alpha) \to U_\alpha\times G, \quad p\mapsto \Big(\pi(p), g_\alpha(p)\Big),
% \end{equation*}
% The definition of $g_\alpha(p)$ is as follows: 
% Let $s_\alpha: U_\alpha\to P$ be a local section defined by $s_\alpha(x) = \phi_\alpha^{-1}(x, e)$. 
% Since the action of $G$ on $P$ is free and transitive on each fiber, for any $p\in \pi^{-1}(U_\alpha)$, 
% there exists a unique $g_\alpha(p)\in G$ such that $p=s_\alpha(\pi(p))g_\alpha(p)$. 
% The transition functions of this local trivialization are given by $g_{\alpha\beta}(x) = g_\alpha(s_\beta(x))$.

% Then we can construct a local trivialization $\{(U_\alpha, \psi_\alpha)\}$ of $E$ by 
% \begin{equation*}
%     \psi_\alpha: \pi_E^{-1}(U_\alpha) \to U_\alpha\times V, \quad [p, v]\mapsto \Big(\pi(p), \rho(g_\alpha(p))v\Big).
% \end{equation*}
% One can check that $\psi_\alpha$ is well-defined and is a diffeomorphism. 
% The transition functions of this local trivialization are given by 
% $\rho\circ g_{\alpha\beta}: U_\alpha\cap U_\beta\to \mathrm{GL}(V)$.

\subsection{Tensorial Forms on a Principal Bundle}
Let $\pi: P\to M$ be a principal $G$-bundle, $\rho:G\to \mathrm{GL}(V)$ be a finite-dimensional representation of $G$, 
and $E=P\times_\rho V$ be the associated vector bundle.
\begin{definition}
    A $V$-valued $k$-form $\varphi\in\Omega^k(P, V)$ is called \textbf{tensorial} if it satisfies the following two conditions:
    \begin{itemize}
        \item $\varphi$ is \textbf{horizontal}: $\varphi(X_1, \cdots, X_k)=0$ if at least one of $X_i$ is vertical.
        \item $\varphi$ is \textbf{$G$-equivariant of type $\rho$}: $r_g^*\varphi = \rho(g^{-1})\cdot\varphi$ for all $g\in G$.
    \end{itemize}
    The set of all tensorial $k$-forms on $P$ of type $\rho$ is denoted by $\Omega^k_{\rho}(P, V)$.
\end{definition}
\begin{example}
    The curvature $\Omega$ of a connection $\omega$ on $P$ is a tensorial 2-form of type $\mathrm{Ad}$.
\end{example}
Below we will show that there is a one-to-one correspondence between $\Omega^k_{\rho}(P, V)$ and $\Omega^k(M, E)$, 
where $\Omega^k(M, E)$ is the space of $E$-valued $k$-forms on $M$. 
\begin{definition}
    Given $x\in M$ and $X_1, \cdots, X_k\in T_xM$, choose any point $p\in P_x$ 
    and choose any lift $\tilde{X}_1, \cdots, \tilde{X}_k\in T_pP$ of $X_1, \cdots, X_k$, i.e., 
    $\tilde{X}_i\in T_pP$ and $\pi_*(\tilde{X}_i) = X_i$. Define $\varphi^\flat$ by
    \begin{equation*}
        \varphi_x^\flat(X_1, \cdots, X_k) := f_p(\varphi_p(\tilde{X}_1, \cdots, \tilde{X}_k)) \in E_x.  
    \end{equation*}
    where $f_p: V\to E_x,\,v\mapsto [p, v]$ is the canonical isomorphism defined above.

    Conversely, given $\psi\in\Omega^k(M, E)$, $p\in P$ and $Y_1, \cdots, Y_k\in T_pP$, define $\psi^\sharp$ by
    \begin{equation*}
        \psi_p^\sharp(Y_1, \cdots, Y_k) := f_p^{-1}\big(\psi_{\pi(p)}(\pi_*Y_1, \cdots, \pi_*Y_k)\big) \in V.
    \end{equation*}
\end{definition}
\begin{theorem}\label{thm:tensorial-forms-correspond-to-forms-on-base}
    These two maps \textcolor{red}{
    \begin{gather*}
        \Omega^k_{\rho}(P, V) \to \Omega^k(M, E), \quad \varphi\mapsto \varphi^\flat,\\
        \Omega^k(M, E) \to \Omega^k_{\rho}(P, V), \quad \psi\mapsto \psi^\sharp,
    \end{gather*}}
    are well-defined and inverse to each other, 
    thus we get linear isomorphisms between $\Omega^k_{\rho}(P, V)$ and $\Omega^k(M, E)$.
\end{theorem}
Before proving the theorem, let's see some examples and corollaries.
\begin{example}
    By the theorem above, the curvature $\Omega$ of a connection on a principal $G$-bundle $P$ 
    can be viewed as an element in $\Omega^2(M,\mathrm{ad}P)$, 
    a $2$-form on $M$ with values in the adjoint bundle $\mathrm{ad}P$.
\end{example}
\begin{corollary}
    Take $k=0$ in Theorem \ref{thm:tensorial-forms-correspond-to-forms-on-base}, 
    an element in $\Omega^0_{\rho}(P, V)$ is just a $G$-equivariant smooth map $f: P\to V$ of type $\rho$, i.e., 
    \begin{equation*}
        f(pg) = (r_g^*f)(p) = \rho(g^{-1})f(p) = g^{-1}\cdot f(p), \quad \forall p\in P, g\in G.
    \end{equation*}
    Besides, $\Omega^0(M, E)$ is just the space of smooth sections of $E$. Hence we get a one-to-one correspondence 
    \begin{equation*}
        \left\{
            \begin{array}{c}
                \text{G-equivariant maps w.r.t. $\rho$} \\[2pt]
                f: P \to V
            \end{array}
        \right\}
        \longleftrightarrow
        \left\{
            \begin{array}{c}
                \text{sections of the associated bundle} \\[2pt]
                s: M \to P \times_\rho V
            \end{array}
        \right\}
    \end{equation*}
\end{corollary}

\begin{proof}[Proof of Theorem \ref{thm:tensorial-forms-correspond-to-forms-on-base}]
    We first show that $\varphi^\flat$ is well-defined. First check that $\varphi^\flat$ does not depend on the choice of the lift $\tilde{X}_i$.
    If $\tilde{X}_i'$ is another lift of $X_i$, then $\pi_*(\tilde{X}_i-\tilde{X}_i')=0$, so $\tilde{X}_i-\tilde{X}_i'$ is vertical. 
    Since $\varphi$ is horizontal, 
    \begin{equation*}
        \varphi_p(\tilde{X}_1', \cdots, \tilde{X}_k') = \varphi_p(\tilde{X}_1+\text{vertical}, \cdots, \tilde{X}_k+\text{vertical}) 
        = \varphi_p(\tilde{X}_1, \cdots, \tilde{X}_k).
    \end{equation*}
    Next check that $\varphi^\flat$ does not depend on the choice of $p\in P_x$. Suppose $p'=pg$ for some $g\in G$, then 
    $r_g{}_*\tilde{X}_i$ is a lift of $X_i$ at $pg$. Since $\varphi$ is $G$-equivariant of type $\rho$, we have
    \begin{equation*}
        \varphi_{pg}(r_g{}_*\tilde{X}_1, \cdots, r_g{}_*\tilde{X}_k) 
        = (r_g^*\varphi_{pg})(\tilde{X}_1, \cdots, \tilde{X}_k)
        = \rho(g^{-1})\cdot\varphi_p(\tilde{X}_1, \cdots, \tilde{X}_k).
    \end{equation*}  
    By Lemma \ref{lem:G-equivariance-of-fp}, 
    \begin{align*}
        f_{pg}(\varphi_{pg}(r_g{}_*\tilde{X}_1, \cdots, r_g{}_*\tilde{X}_k)) 
        = f_p\circ \rho(g)(\rho(g^{-1})\cdot\varphi_p(\tilde{X}_1, \cdots, \tilde{X}_k))
        = f_p(\varphi_p(\tilde{X}_1, \cdots, \tilde{X}_k)).
    \end{align*}
    Thus $\varphi^\flat$ is well-defined. 

    Let $\psi\in\Omega^k(M, E)$, it's easy to see that $\psi^\sharp$ is horizontal. For $G$-equivariance, we have 
    \begin{align*}
        (r_g^*\psi^\sharp)_p(Y_1, \cdots, Y_k) 
        &= \psi^\sharp_{pg}(r_g{}_*(Y_1), \cdots, r_g{}_*(Y_k))\\
        &= f_{pg}^{-1}\big(\psi_{\pi(pg)}(\pi_*r_g{}_*(Y_1), \cdots, \pi_*r_g{}_*(Y_k))\big)\\
        &= f_{pg}^{-1}\big(\psi_{\pi(p)}(\pi_*Y_1, \cdots, \pi_*Y_k)\big)\\
        &= \rho(g^{-1})\cdot f_p^{-1}\big(\psi_{\pi(p)}(\pi_*Y_1, \cdots, \pi_*Y_k)\big)\\
        &= \rho(g^{-1})\cdot \psi^\sharp_p(Y_1, \cdots, Y_k).
    \end{align*}
    Hence $\psi^\sharp$ is $G$-equivariant of type $\rho$, and thus $\psi^\sharp\in\Omega^k_{\rho}(P, V)$.

    Finally, we check that these two maps are inverse to each other. For any $\varphi\in\Omega^k_{\rho}(P, V)$, 
    $p\in P$ and $Y_1, \cdots, Y_k\in T_pP$, 
    \begin{align*}
        (\varphi^{\flat\sharp})_p(Y_1, \cdots, Y_k) 
        &= f_p^{-1}\big(\varphi^\flat_{\pi(p)}(\pi_*Y_1, \cdots, \pi_*Y_k)\big)\\
        &= f_p^{-1}\big(f_p(\varphi_p(Y_1, \cdots, Y_k))\big)\\
        &= \varphi_p(Y_1, \cdots, Y_k).
    \end{align*}
    For any $\psi\in\Omega^k(M, E)$, $x\in M$ and $X_1, \cdots, X_k\in T_xM$, 
    choose any $p\in P_x$ and any lift $\tilde{X}_1, \cdots, \tilde{X}_k\in T_pP$ of $X_1, \cdots, X_k$, 
    \begin{align*}
        (\psi^{\sharp\flat})_x(X_1, \cdots, X_k) 
        &= f_p(\psi^\sharp_p(\tilde{X}_1, \cdots, \tilde{X}_k))\\
        &= f_p\big(f_p^{-1}(\psi_{x}(\pi_*\tilde{X}_1, \cdots, \pi_*\tilde{X}_k))\big)\\
        &= \psi_{x}(X_1, \cdots, X_k). \qedhere
    \end{align*}
\end{proof}
\subsection{Covariant Derivative}
Keep the same notations as above. Notice that given a horizontal distribution $H$ on $P$, 
we can decompose the tangent vetor $Y_p\in T_pP$ into a horizontal part and a vertical part, 
\begin{equation*}
    Y_p = hY_p + vY_p, \quad hY_p\in H_p, vY_p\in V_pP.
\end{equation*}
In this section, we will use this decomposition to define the covariant derivative of vector-valued forms on $P$, 
and thus a covariant derivative on the associated bundle $E$. 

Let's start with the usual exterior derivative $\md: \Omega^k(P, V)\to \Omega^{k+1}(P, V)$. 
\begin{proposition}\label{prop:exterior-derivative-preserves-G-equivariance}
    If $\varphi\in\Omega^k(P, V)$ is $G$-equivariant of type $\rho$, then $\md\varphi$ is also $G$-equivariant of type $\rho$.
\end{proposition}
\begin{proof}
    For any $g\in G$, 
    \begin{equation*}
        r_g^*(\md\varphi) = \md(r_g^*\varphi) = \md(\rho(g^{-1})\varphi) = \rho(g^{-1})\md\varphi.
    \end{equation*}
    The last equality holds since $\rho(g^{-1})$ is a constant linear map on $V$ for a fixed $g$.
\end{proof}

In general, the exterior derivative does not preserve the horizontality, so it does not preserve the space of tensorial forms.
\begin{definition}
    For any $\varphi\in\Omega^k(P, V)$, we define its \textbf{horizontal component} $\varphi^h\in\Omega^k(P, V)$ by
    \begin{equation*}
        \varphi^h_p(Y_1, \cdots, Y_k) := \varphi_p(hY_1, \cdots, hY_k),
    \end{equation*}
    for any $p\in P$ and $Y_1, \cdots, Y_k\in T_pP$. 
\end{definition}
\begin{proposition}\label{prop:horizontal-component-preserves-G-equivariance}
    If $\varphi\in\Omega^k(P, V)$ is $G$-equivariant of type $\rho$, so is $\varphi^h$.
\end{proposition}
\begin{proof}
    For $g\in G$, $p\in P$ and $Y_1, \cdots, Y_k\in T_pP$,
    \begin{align*}
        (r_g^*\varphi^h)_p(Y_1, \cdots, Y_k) 
        &= \varphi^h_{pg}(r_g{}_*Y_1, \cdots, r_g{}_*Y_k)\\
        &= \varphi_{pg}(h(r_g{}_*Y_1), \cdots, h(r_g{}_*Y_k))\\
        &= \varphi_{pg}(r_g{}_*(hY_1), \cdots, r_g{}_*(hY_k))\\
        &= (r_g^*\varphi)_p(hY_1, \cdots, hY_k)\\
        &= \rho(g^{-1})\cdot\varphi_p(hY_1, \cdots, hY_k)\\
        &= \rho(g^{-1})\cdot\varphi^h_p(Y_1, \cdots, Y_k).
    \end{align*}
\end{proof}

Proposition \ref{prop:exterior-derivative-preserves-G-equivariance} and \ref{prop:horizontal-component-preserves-G-equivariance} 
together imply that if $\varphi\in\Omega^k(P, V)$ is $G$-equivariant of type $\rho$, then $(\md\varphi)^h$ is a tensorial form.
\begin{definition}
    For any $\varphi\in\Omega^k(P, V)$, we define its \textbf{covariant derivative} by 
    \begin{equation*}
        D\varphi := (\md\varphi)^h.
    \end{equation*}
    In particular, $D$ maps tensorial $k$-forms of type $\rho$ to tensorial $(k+1)$-forms of type $\rho$, i.e., 
    \begin{equation*}
        D: \Omega^k_{\rho}(P, V) \to \Omega^{k+1}_{\rho}(P, V).
    \end{equation*}
\end{definition}
\begin{proposition}
    The covariant derivative $D: \Omega^k_{\rho}(P, V) \to \Omega^{k+1}_{\rho}(P, V)$ is an antiderivation of degree $+1$, 
    i.e., for any $\varphi\in\Omega^k_{\rho}(P, V)$ and $\psi\in\Omega^l_{\rho}(P, V)$,
    \begin{equation*}
        D(\varphi\wedge\psi) = D\varphi\wedge\psi + (-1)^k\varphi\wedge D\psi.
    \end{equation*}
\end{proposition}
\begin{proof}
    Since $\varphi$ and $\psi$ are horizontal, we have 
    \begin{align*}
        D(\varphi\wedge\psi) &= (\md(\varphi\wedge\psi))^h = (\md\varphi\wedge\psi + (-1)^k\varphi\wedge \md\psi)^h  \\ 
        &= (\md\varphi)^h\wedge\psi^h + (-1)^k\varphi^h\wedge (\md\psi)^h \\ 
        &= D\varphi\wedge\psi + (-1)^k\varphi\wedge D\psi. \qedhere
    \end{align*}
\end{proof}
\begin{example}
    By Theorem \ref{thm:tensorial-forms-correspond-to-forms-on-base}, 
    the curvature $\Omega$ of a connection on $P$ is exactly the covariant derivative of $\omega$, since
    \begin{equation*}
        \Omega_p(X_p, Y_p) = \md\omega_p(hX_p, hY_p) = (\md\omega_p)^h(X_p, Y_p) = (D\omega_p)(X_p, Y_p). 
    \end{equation*}
\end{example}
By the one-to-one correspondence between $\Omega^k_{\rho}(P, V)$ and $\Omega^k(M, E)$, we can also define a covariant derivative on $\Omega^k(M, E)$ 
\begin{equation*}
    D: \Omega^k(M, E) \to \Omega^{k+1}(M, E), \quad D\psi := (D\psi^\sharp)^\flat.
\end{equation*}
\begin{example}
    Let $E$ be a rank $m$ vector bundle over $M$, $\mathrm{Fr}(E)$ be its frame bundle.
    and $\omega$ be a connection on $\mathrm{Fr}(E)$ induced by a connection $\nabla$ on $E$. 
    With the identification $E = \mathrm{Fr}(E)\times_{\rho_{\mathrm{can}}}\mathbb{R}^m$, 
    we claim that the covariant derivative $D$ on $\Omega^k(M, E)$ is exactly $\nabla$. 
    The proof will be given later. 
\end{example}

\subsection{A Formula for the Covariant Derivative on Tensorial Forms}
We have defined the covariant derivative $D$ on $\Omega^k(P, V)$ by $D\varphi = (\md\varphi)^h$.
In this section, we will give a more explicit formula for $D$, but only for $\Omega^k_{\rho}(P, V)$. 

The representation $\rho:G\to \mathrm{GL}(V)$ induces a Lie algebra representation $\rho_*:\mathfrak{g}\to \mathrm{End}(V)$, 
which alows us to define a product of $\tau\in\Omega^k(P,\mathfrak{g})$ and $\varphi\in\Omega^l(P, V)$ 
\begin{definition}
    For $\tau\in\Omega^k(P,\mathfrak{g})$ and $\varphi\in\Omega^l(P, V)$, we define $\tau\cdot\varphi\in\Omega^{k+l}(P, V)$ by
    \begin{equation*}
        (\tau\cdot\varphi)_p(Y_1, \cdots, Y_{k+l}) := \frac{1}{k!l!}\sum_{\sigma} \mathrm{sgn}(\sigma)\cdot 
        \rho_*\big(\tau_p(Y_{\sigma(1)}, \cdots, Y_{\sigma(k)})\big)\varphi_p(Y_{\sigma(k+1)}, \cdots, Y_{\sigma(k+l)}).
    \end{equation*}
    One can check that $\tau\cdot\varphi$ is indeed a $(k+l)$-form on $P$ with values in $V$.
\end{definition}
\begin{example}
    If $V=\mathfrak{g}$ and $\rho=\mathrm{Ad}$, then 
    \begin{equation*}
        (\tau\cdot\varphi)_p(Y_1, \cdots, Y_{k+l}) = \frac{1}{k!l!}\sum_{\sigma} \mathrm{sgn}(\sigma)\cdot 
        [\tau_p(Y_{\sigma(1)}, \cdots, Y_{\sigma(k)}), \varphi_p(Y_{\sigma(k+1)}, \cdots, Y_{\sigma(k+l)})].
    \end{equation*}
    In this case, we write $\tau\cdot\varphi$ instead of $[\tau, \varphi]$.
\end{example}
\begin{theorem}\label{thm:formula-for-covariant-derivative}
    Let $\pi: P\to M$ be a principal $G$-bundle with a connection $\omega$, 
    $\rho:G\to \mathrm{GL}(V)$ be a finite-dimensional representation of $G$. 
    If $\varphi\in\Omega^k_{\rho}(P, V)$, then 
    \begin{equation*}
        D\varphi = \md\varphi + \omega\cdot\varphi.
    \end{equation*}
\end{theorem}
Before proving the theorem, we frist see some examples. 
\begin{example}
    In case that $V=\mathfrak{g}$ and $\rho=\mathrm{Ad}$, for any $\varphi\in\Omega^k_{\mathrm{Ad}}(P, \mathfrak{g})$, 
    \begin{equation*}
        D\varphi = \md\varphi + [\omega, \varphi].
    \end{equation*}
    Notice that this formula only holds for tensorial forms. Since $\omega$ is not horizontal, 
    we cannot use this formula to compute $D\omega$, in fact, we know that 
    \begin{equation*}
        D\omega = (\md\omega)^h = \Omega = \md\omega + \frac{1}{2}[\omega, \omega].
    \end{equation*}
    Since $\Omega$ is tensorial, we can apply the formula for $\Omega$. In fact, we have
    \begin{equation*}
        D\Omega = \md\Omega + [\omega, \Omega].
    \end{equation*}
    So the second Bianchi identity can be rewritten as $D\Omega=0$.
\end{example}
\begin{proof}[Proof of Theorem \ref{thm:formula-for-covariant-derivative}]
    To verify the formula, we actually need to show
    \begin{equation*}
        \md\varphi(hY_1, \cdots, hY_{k+1}) = \md\varphi(Y_1, \cdots, Y_{k+1}) + (\omega\cdot\varphi)_p(Y_1, \cdots, Y_{k+1}) 
        =: \mathrm{I} + \mathrm{II}.
    \end{equation*}
    Since both sides are multilinear and antisymmetric in $Y_1, \cdots, Y_{k+1}$, 
    we can decompose all the $Y_i$ into horizontal and vertical parts and check the equality in several cases. 
    Moreover, since we are verifying the equalityly, we can assume that the vertical part of $Y_i=\underline{A_i}$ is a fundamental vector field, 
    and the horizontal part of $Y_i=\tilde{B}_i$ is a horizontal lift of some vector field $B_i$ on $M$. 
    Notice that $\tilde{B}_i$ is right-invariant. 
    
    \noindent\textbf{Case 1:} All $Y_i$ are horizontal. Then $\mathrm{II}=0$ since $\omega$ vanishes on horizontal vectors, 
    and the formula holds trivially.

    \noindent\textbf{Case 2:} At least two of the $Y_i$ are vertical. Then 
    \begin{align*}
        \mathrm{I} = \md\varphi(Y_1, \cdots, Y_{k+1}) 
        =& \sum_i (-1)^{i+1}Y_i\varphi(Y_1, \cdots, \hat{Y}_i, \cdots, Y_{k+1}) \\ 
        &+ \sum_{i<j} (-1)^{i+j}\varphi([Y_i, Y_j], Y_1, \cdots, \hat{Y}_i, \cdots, \hat{Y}_j, \cdots, Y_{k+1}).
    \end{align*}
    Since $[\underline{A_i}, \underline{A_j}] = \underline{[A_i, A_j]}$ is still vertical, 
    at least one of the arguments is vertical, so $\mathrm{I}=0$. 

    On the other hand, 
    \begin{equation*}
        \mathrm{II} = \frac{1}{k!}\sum_{\sigma\in S_{k+1}} \mathrm{sgn}(\sigma)\cdot \rho_*(\omega_p(Y_{\sigma(1)}))\varphi_p(Y_{\sigma(2)}, \cdots, Y_{\sigma(k+1)}),
    \end{equation*}
    which is also zero since at least one of the arguments of $\varphi$ is vertical. 

    \noindent\textbf{Case 3:} Exactly one of the $Y_i$ is vertical, say $Y_1=\underline{A}$ and $Y_2, \cdots, Y_{k+1}$ are horizontal. 
    Then 
    \begin{equation*}
        \mathrm{I} = \md\varphi(\underline{A}, Y_2, \cdots, Y_{k+1}) 
        = \underline{A}\varphi(Y_2, \cdots, Y_{k+1}) + \sum_{i=2}^{k+1} (-1)^{i}\varphi([\underline{A}, Y_i], Y_2, \cdots, \hat{Y}_i, \cdots, Y_{k+1}).
    \end{equation*}
    In our assumption, $Y_i$ is a horizontal lift of some vector field $B_i$ on $M$, so is right-invariant. 
    Hence $[\underline{A}, Y_i] = 0$ for $i=2, \cdots, k+1$. Thus 
    \begin{equation*}
        \mathrm{I} = \underline{A}\varphi(Y_2, \cdots, Y_{k+1}).
    \end{equation*}
    On the other hand,
    \begin{align*}
        \mathrm{II} &= \frac{1}{k!}\sum_{\begin{subarray}{c}\sigma\in S_{k+1} \\ \sigma(1)=1\end{subarray}}
        \mathrm{sgn}(\sigma)\cdot \rho_*(\omega(Y_{\sigma(1)}))\varphi(Y_{\sigma(2)}, \cdots, Y_{\sigma(k+1)}) \\ 
        &= \rho_*(\omega(\underline{A}))\varphi(Y_2, \cdots, Y_{k+1})  
        = \rho_*(A)\varphi(Y_2, \cdots, Y_{k+1}).
    \end{align*}
    If we denote by $f$ the function $\varphi(Y_2, \cdots, Y_{k+1})$. For $p\in P$, to calculate $\underline{A}_pf$, 
    let $c(t) = \exp(tA)$ be the integral curve of $A$ in $G$ with $c(0)=e$. Then withe $j_p: G\to P, j_p(g) = pg$, we have
    \begin{align*}
        \underline{A}_pf = j_p{}_*(A)f = j_p{}_*(c'(0))f = \frac{\md}{\md t}\Big|_{t=0} (f\circ j_p\circ c).
    \end{align*}
    By the right-invariance of the horizontal vector fields $Y_2, \cdots, Y_{k+1}$, 
    \begin{align*}
        (f\circ j_p\circ c)(t) &= f(pc(t)) = \varphi_{pc(t)}(Y_{2,pc(t)}, \cdots, Y_{k+1, pc(t)}) \\ 
        &= \varphi_{pc(t)}(r_{c(t)*}Y_{2,p}, \cdots, r_{c(t)*}Y_{k+1, p}) \\
        &= (r_{c(t)}^*\varphi)_p(Y_{2,p}, \cdots, Y_{k+1, p}) \\
        &= \rho(c(t)^{-1})\varphi_p(Y_{2,p}, \cdots, Y_{k+1, p}) \\
        &= \rho(\exp(-tA))f(p).
    \end{align*}
    Thus 
    \begin{equation*}
        \underline{A}_pf(p) = \frac{\md}{\md t}\Big|_{t=0} \rho(\exp(-tA))f(p) = -\rho_*(A)f(p).
    \end{equation*}
    Hence $\mathrm{I} + \mathrm{II} = 0$. The left-hand side of the formula is also zero because $hY_1 = 0$. 
\end{proof}
\subsection*{Exercises}
\addcontentsline{toc}{subsection}{Exercises}
\begin{exercise}[Forms on $P$ pulled back from $M$]
    Since $\pi:P\to M$ is a surjective submersion, 
    the pull-back $\pi^*:\Omega^*(M,V)\to \Omega^*(P,V)$ is injective. 
    A differential form $\varphi\in\Omega^*(P,V)$ is said to be \textbf{basic} 
    if it is the pull-back $\pi^*\psi$ of some $\psi\in\Omega^*(M,V)$; 
    it is said to be \textbf{$G$-invariant} if $r_g^*\varphi = \varphi$ for all $g\in G$. 
    Show that:
    \begin{itemize}[leftmargin=*]
        \item $\varphi$ is $G$-invariant if and only if 
        $\varphi$ is $G$-equivariant of type $\rho_{\mathrm{trivial}}$, 
        where $\rho_{\mathrm{trivial}}$ is the trivial representation.
        \item $\varphi$ is basic if and only if $\varphi$ is horizontal and $G$-invariant.
    \end{itemize}
\end{exercise}
\begin{proof}
    If $\varphi$ is $G$-equivariant of type $\rho_{\mathrm{trivial}}$, then 
    \begin{equation*}
        r_g^*\varphi = \rho_{\mathrm{trivial}}(g^{-1})\cdot\varphi = \varphi, \quad \forall g\in G.
    \end{equation*}
    This is just the definition of $G$-invariance. 

    Since the associated bundle $E=P\times_{\rho_{\mathrm{trivial}}} V$ is just the trivial bundle $M\times V$, 
    and $f_p: V\to E_{\pi(p)} = V$ is the identity map,
    by Theorem \ref{thm:tensorial-forms-correspond-to-forms-on-base}, 
    \begin{align*}
        \Omega^k(M,E)\cong\Omega^k(M, M\times V) = \Omega^k(M, V) &\xrightarrow{\cong} \Omega^k_{\rho_{\mathrm{trivial}}}(P, V) = \{\text{basic forms on $P$}\} \\ 
        \psi &\mapsto \psi^\sharp.
    \end{align*}
    where 
    \begin{align*}
        \psi^\sharp_p(Y_1, \cdots, Y_k) &= f_p^{-1}\Big(\psi_{\pi(p)}(\pi_*Y_1, \cdots, \pi_*Y_k)\Big) \\ 
        &= \psi_{\pi(p)}(\pi_*Y_1, \cdots, \pi_*Y_k) \\ 
        &= (\pi^*\psi)_p(Y_1, \cdots, Y_k). \qedhere
    \end{align*}
\end{proof}
